The measured interval is the host search call:
\[
 t_{\rm call}=t_{\rm return}-t_{\rm entry},
 \quad
 \begin{cases}
 \text{entry: FP32 queries in host memory},\\
 \text{return: host IDs/scores after stream completion}.
 \end{cases}
\]
It includes staging, transfers, query work and int8 query quantisation; fitting/build are excluded (\cref{fig:3-boundary}).
Resource handles and output buffers are reused. Output/API differences remain: JHQ pageable/C++; CAGRA/PQ pinned/Python; RaBitQ pinned/C++. Small margins include these costs.

\par\medskip\noindent\begin{minipage}{\linewidth}
\centering\captionsetup{hypcap=false}
\resizebox{\linewidth}{!}{% Protocol diagram, not a measured stage-time plot.
\begin{tikzpicture}[
 font=\small\linespread{1}\selectfont,
 box/.style={draw=ThesisNavy!65,rounded corners=2pt,align=center,
 text width=2.55cm,minimum height=1.55cm,inner sep=4pt},
 arr/.style={-{Latex[length=2mm]},draw=ThesisNavy!80,thick}]
\node[box] (host) at (0,0) {Host FP32 queries\\Pinned staging};
\node[box,fill=ThesisNavy!7] (upload) at (3.5,0) {H2D transfer\\Query preprocessing};
\node[box,fill=ThesisNavy!7] (gpu) at (7,0) {GPU search\\Selection / refinement};
\node[box] (result) at (10.5,0) {D2H IDs and scores\\Stream completion};
\draw[arr] (host) -- (upload);
\draw[arr] (upload) -- (gpu);
\draw[arr] (gpu) -- (result);
\draw[ThesisNavy,thick] (-1.4,1.15) -- (-1.4,1.45) -- (11.9,1.45) -- (11.9,1.15);
\node[fill=white,inner sep=3pt] at (5.25,1.45) {Primary host-to-host timed interval};
\node[box,text width=5.7cm,minimum height=1.3cm] (setup) at (1.65,-2.35)
 {Outside search timing\\Load data / train or load model / add / warm-up};
\node[box,text width=5.7cm,minimum height=1.3cm] (cal) at (8.65,-2.35)
 {Separate policy cost\\Calibration; reuse cost model};
\draw[densely dashed,ThesisNavy!60] (-1.4,-1.1) -- (11.9,-1.1);
\end{tikzpicture}
}
\captionof{figure}[Host and device timing boundaries]{Timing boundaries. The host interval includes transfers and
query-specific computation. Device-stage profiles explain parts of that
interval; separate profiling and synchronisation can change execution, so
their sum is not an interchangeable end-to-end latency.}
\label{fig:3-boundary}
\end{minipage}\par\medskip

With $t$ in seconds,
\begin{equation}
 \bar t_r=\frac1{m_r}\sum_{i=1}^{m_r}t_{ri},
 \qquad Q_r=\frac b{\bar t_r},
 \qquad Q=\operatorname{median}_r Q_r.
 \label{eq:3-throughput}
\end{equation}
E1 uses batch 1,000, five rounds and normally 24 calls. JHQ warms once per variant/round; baselines warm three times initially. Slow baseline points use fewer calls. Warm-up excludes JIT/capture; percentiles describe batch calls, not concurrent requests.

Slow valid calls stay. Mean-call and median-call rates differ; for an illustrative $(1,1,10)$\,ms sequence,
\[
 \bar t=4\,\mathrm{ms}\ne\operatorname{median}(t)=1\,\mathrm{ms},
 \qquad b/\bar t\ne b/\operatorname{median}(t).
\]
Shared indexes/pools are not independent training replications.

Runs use RTX 5090/32,607\,MiB, driver 595.71.05, CUDA 13.0 and cuVS 26.08.01. OpenMP uses 16 threads within 25 CPUs; E1 records 575\,W. Locks do not fix clocks/load. JHQ JIT-compiles compute\_80 PTX \cite{cuda_platform}; library builds and device profiles remain separate.

Resident totals include models, workspace and context; CAGRA's index deltas differ from JHQ's context-subtracted totals (\cref{sec:4-4,sec:7-7}). Later pools and build scratch prevent peak-matched comparison. Build peaks use 20\,ms sampling; train/add excludes input loading and host setup.

\paragraph{Supplementary evidence boundary.}
Supplementary sessions are separate (\cref{sec:a-supplement}). Precision changes and cached-model reuse follow \cref{sec:4-2,sec:7-precision}; historical baseline ratios remain cross-session.